\section{Families from iterated theta lifts: the hypothesis package}\label{sec:families}

\providecommand{\familiesstr}{\operatorname{str}}
\providecommand{\familiesClosed}{\sfname{Closed}}
\providecommand{\familiesLinked}{\sfname{LinkedAt}}
\providecommand{\familiesRow}{\sfname{Row}}
\providecommand{\familiesDiii}{\mathcal{D}_3}

In \cref{sec:top} the main theorem was reduced to three propositions: Green's theorem, the window bound
$\sfname{VWin}$, and the non-emptiness of the hypothesis package $\sfname{HypsAlt}$ of \cref{def:hyps}, with the
parameters $\sfname{NL}$ and $\sfname{DefiniteSideCondition}$. In
\cref{sec:stepanov}
we explained how an element of the package, together with Green's theorem, yields the proposition
$\sfname{TopG}$ (\cref{thm:shape}), and how the chain-tied records behind $\sfname{TopG}$ are obtained from a
list of \emph{links} by the global theorem (\cref{def:links,thm:global-chain}). The present section
explains where the package comes from and how the individual links are established. Its main result is
\cref{thm:hyps-nonempty}: the package is non-empty, given two bundles of classical results on unitary groups
(all of them proved in the formal development), the theta-integral identities of Theorem~K
(\cref{thm:theorem-K}), the covolume bound of \cref{sec:covolume} and the local parameters of
\cref{sec:galois}.

The objects involved are large records, and the formal statements are inequalities between their components.
We name the components that matter, state the results as they are proved, and outline the proofs; complete
proofs are those of the formal development \cite{LeanSource2026} (results marked~\Fmark), verified as described in
\cite{Sutherland2026}. The picture of \cite{Announcement2026} recalled below is a guide to the meaning of the
formal objects; we claim only what the formal statements give.

% ---------------------------------------------------------------------------
\subsection{Overview}\label{sec:families:overview}

\paragraph{The picture.}
Stepanov's method bounds the number of rational points of a curve over $\F_q$ by means of an auxiliary function
vanishing to order $m\asymp\sqrt q$ at every rational point. In the present argument the multiplicity is
supplied by chains of theta lifts
\begin{equation}\label{eq:families:chain}
\Uni(X_0)\longrightarrow\Uni(X_1)\longrightarrow\cdots\longrightarrow\Uni(X_m)
\end{equation}
between totally definite unitary groups attached to the CM host field $E$ of \cref{sec:hosts}
\cite{Announcement2026}. On unramified Satake parameters the $j$-th lift adds a segment of length $t$
twisted by a Hecke character $\mu_j$ (in the notation of Arthur parameters, $\pi\mapsto\pi\boxplus\mu_j[t]$;
see \cref{def:families:string}), and each lift plays the role of one order of vanishing. Accordingly the
length $m$ of the chain is of the order of the square root of the scale: the rank records admitted by the
host-existence condition $\sfname{HostExW}$ (\cref{def:hostexw}) satisfy
\begin{equation}\label{eq:families:rankval}
x^{2+8\eps}\le n\le 2x^{2+8\eps},\qquad (x\delta_1)^{1/2}\le m\le 2(x\delta_1)^{1/2},\qquad
x^{9/2+10\eps}\le k\le 2x^{9/2+10\eps},
\end{equation}
with $\delta_1=x^{-\eps/8}$ (condition $\sfname{RankVal}$ of the development), so that $m$ is of size
$x^{1/2-\eps/16}$.

From the lifts one builds integral structures: hermitian lattices over the ring of integers of a coefficient
field $\mathit{LL}$, indexed by characters $\mu$ and by ``words'' $u$, together with a lattice
$\mathit{RR}$ glued from local lattices. Their local indices at a finite place $w$ measure the $w$-adic
divisibility of the iterated lifts; their archimedean size is measured by covolumes and by slopes in the sense
of Stuhler and Grayson \cite{Stuhler1976,Grayson1984}. Three groups of results control these quantities.
\begin{enumerate}[label=(\alph*)]
\item \emph{Local divisibility.} At the places above the primes in the inert Frobenius classes of the host,
the gains in divisibility of the lattices dominate the von Mangoldt sums over these classes, and at the level
primes the lattices are divisible to an order proportional to $m$. The inputs are theta functions of CM
elliptic curves, whose periods are governed by the Chowla--Selberg formula \cite{ChowlaSelberg1967}, lattices in
Weyl modules \cite{Weyl1946}, and a local lattice theorem (\cref{thm:linkG,thm:families:local}).
\item \emph{Archimedean size.} Covolumes, masses and norms of lifts are computed through the Siegel--Weil
formula, the Siegel mass formula for hermitian lattices and the Rallis inner product formula, and enter slope
inequalities of the kind used by Bost \cite{Bost1996,Bost2001} (\cref{sec:families:rows,sec:families:links,sec:families:nonempty}).
\item \emph{Spectral and Galois-theoretic control.} Theta series attached to distinct words are approximately
orthogonal (Theorem~S, \cref{thm:theorem-S}); spaces of automorphic forms are counted with exponentially small
relative error (\cref{thm:counting}); Green's theorem controls the irreducible constituents
(\cref{thm:linkIRR}); and the failure of divisibility at the twisted places is bounded by lengths of Selmer
groups and of $\Sha^1(G_{E,S_\ell},\Ql/\Zl(-e))$ (\cref{thm:linkN}).
\end{enumerate}
These inequalities are combined by the global theorem of \cref{sec:stepanov} (\cref{thm:global-chain}).

\paragraph{What must be supplied.}
\Cref{tab:families:links} lists the links of \cref{def:links} and the results of this section that establish
them. All links are proved as \emph{closed} properties, i.e.\ for all families of a class with small exponent
$\eps$ and sufficiently large constants (\cref{def:families:closed}), on the class
$\familiesDiii=D3(3,\mathcal A)$ attached to the class $\mathcal A$ of an element of the package; the
hypotheses of the link theorems are fields of the package. Throughout, $D=D_{x,i}$ and $L=L_{x,i}$ denote the
datum and the size data of a member of a family, $b=\exp(x^\eps)$ and $\Lambda=x^\eps$ are the scale quantities
of $L$, and $j$ ranges over the lift indices $1\le j\le m$.

{\small
\begin{longtable}{>{\raggedright\arraybackslash}p{0.2\textwidth}>{\raggedright\arraybackslash}p{0.46\textwidth}>{\raggedright\arraybackslash}p{0.24\textwidth}}
\caption{The links and where they are established.}\label{tab:families:links}\\
\toprule
link & content (eventually in $x$, for all $i\in I_x$) & established in\\
\midrule
\endfirsthead
\toprule
link & content & established in\\
\midrule
\endhead
$\sfname{LinkN}$ & weighted divisibility defect $\eta\Sigma(L,j)$ at twisted places bounded by
$d_{Z_0}(j)\bigl(C(fmb\Lambda^2+(\log x)^2N_c)+\sum_\ell\log\ell\sum_{0<e<t}\mathrm{sh}_\ell(E,e)\bigr)$
& \cref{thm:linkN}; field $\mathsf{nRest}$\\
$\sfname{LinkIRR}$ & $0\le\delta_p(D)\le\exp(-x^{\eps/(7J)})$: almost all labelled indices are rows
& \cref{thm:linkIRR} (Green)\\
$\sfname{LinkG}$ & local lattice lengths at places $w$; at places over primes $p\le X''(n)$ in the inert
classes, $0\le$ gains $\le4|TT|a/p$, and the gains dominate the von Mangoldt class sums at $p$
& \cref{thm:linkG,thm:families:local}\\
$\sfname{LinkS}$ & Hilbert--Schmidt approximate orthogonality of the spaces attached to distinct words
& \cref{thm:theorem-S,cor:families:linkS}\\
$\sfname{LinkDIVt}$ & under $\sfname{LinkEta}$ and $\sfname{DIVinN}$: frame, height and Bessel identities for
the division lattices & \cref{prop:families:div}\\
$\sfname{LinkEta}$, $\sfname{SmallLevelRankEq}$ & local places cover the twisted places;
$d_{Z_0}(j)=D_0$ & \cref{rem:families:small}\\
$\sfname{NuFullF}$, $\sfname{MuRank}$, $\sfname{LinkGLOBAL}$ & mass ratios, volume $=$ mass in rank $r$, slope
inequalities, column valuations & \cref{prop:families:famgl}\\
design & existence of large designed families & \cref{prop:families:design}\\
\bottomrule
\end{longtable}
}

\paragraph{How the package is built.}
An element of $\sfname{HypsAlt}(\mathcal N,\mathrm{DR})$ consists of a family class $\mathcal A$, a
\emph{selection} $\sigma$ of ``bundles'' for its members, and $44$ further fields (\cref{def:hyps}). The
construction has three layers: the theta correspondence (\cref{sec:families:theta}); \emph{rows}, i.e.\
properties of all bundles linked to the points of a construction, proved eventually in the scale and for large
twisting data (\cref{sec:families:rows,sec:families:galois,sec:families:links}); and the existence of members and
linked bundles, from which a class and a selection satisfying all rows are extracted
(\cref{sec:families:nonempty}).

% ---------------------------------------------------------------------------
\subsection{Unitary groups, theta correspondence and iterated lifts}\label{sec:families:theta}

\paragraph{Global data.}
The records of \cref{def:record} are built on \emph{global data} $D$. Their components used here are:
\emph{ranks} $(k,m,t,r,f,n,\dots)$ with $k,m$ even, $t$ odd, $2\le m$, $m+3\le k$, $k+m+4\le t$, $3m<t$,
$m\le f\le k$ and $m^2\le k+m(t+1)$ ($m$ is the length of the chain \eqref{eq:families:chain} and $t$ the length
of the segments); a \emph{host}: a CM extension $E/F$, distinguished primes $q,d$ and \emph{type primes}
$g_0,g_1\ge5$, a finite set of small primes containing $2$, and the imaginary quadratic field $K=\Q(\sqrt{-d})$;
a totally complex \emph{coefficient field} $\mathit{LL}$, Galois over $\Q$, with characters
$\mu_w\in\{\mu_0,\mu^{\mathrm{tw}}_w\}$ at its finite places ($w$ is \emph{twisted} if $\mu_w\ne\mu_0$); hermitian
spaces over $E$ ($Y_j$, $X'$, and $V_i,W_i$ for $0\le i<m$) with Weil representations and theta data; a finite
set of \emph{words} $u$ with Harish-Chandra parameters and level primes $T_u=\pi\bar\pi$; and Galois data
(\cref{sec:families:galois}). Size data $L$ consist of a scale $(x,\eps,J)$ with $x>1$, $0<\eps<1/52$, $J\ge1$, a
finite support of places of $\mathit{LL}$, local weights, and derived quantities such as $\Lambda=x^\eps$,
$b=e^{x^\eps}$, $N_c$ and the dimensions $d_{Z_0}(j)$.

\paragraph{The two bundles of classical results.}
The construction starts from two propositions,
$P:\sfname{UnitaryGroupThirdPlaceHypotheses}$ and $F:\sfname{ThetaLocalInputs}$, conjunctions of $60$ and $53$
classical results: Harish-Chandra's theory of discrete series of $\Uni(p,q)$
\cite{HarishChandra1966,HarishChandra1968}, Casselman--Wallach globalisation, Dixmier--Malliavin and Flath
factorisation, rapid decay and trace class for cusp forms; Weil representations of unitary dual pairs
\cite{Kudla1994}, Howe duality at the real place, the Siegel--Weil and Weil--Siegel formulas in the convergent
range \cite{Weil1965,KudlaRallis1988,Ichino2004}, the doubling method \cite{GelbartPiatetskiShapiroRallis1987} and
the Rallis inner product formula \cite{Rallis1984}; Bernstein--Zelevinsky theory and Green's cuspidal characters
\cite{Green1955}; class field theory, CM elliptic curves and Chebotarev's theorem; local Tate duality, Fontaine's
$B_{\mathrm{cris}}$, the Bloch--Kato conditions and the comparison of Bhatt--Scholze; Galois representations of
polarized cuspidal representations, base change for definite unitary groups, depth-zero support and partial
Hecke $L$-functions. Every field of $P$ and $F$ is a theorem of the formal development (\cref{sec:literature});
four fields of $P$ are deduced from the statements of \cref{sec:galois}. (One field of $P$, named after the simple
pole of Siegel Eisenstein series, is the trivial proposition; the pole used in \cref{prop:families:liftinj} is
obtained from partial Hecke $L$-functions.)

\paragraph{Complementary spaces and Theorem K.}
Let $E/F$ be a CM extension, $V$ and $V'$ hermitian spaces of dimensions $a>a'$, and $W$ a skew-hermitian
space of dimension $a+a'$ (which is even). Let $\alpha,\beta$ be Hecke characters of $\A_E^\times$, let $\omega$
be a Weil representation of $\Uni(V)\times\Uni(W)$ with splitting characters $(\beta,\beta/\alpha)$ and
$\omega'$ one of $\Uni(V')\times\Uni(W)$ with splitting characters $(\alpha,\alpha/\beta)$, and consider the
theta integrals
\[
I(\Phi)(h)=\int_{\Uni(V)(F)\backslash\Uni(V)(\A)}\theta(g,h;\Phi)\,dg,\qquad
I'(\Phi')(h)=\int_{\Uni(V')(F)\backslash\Uni(V')(\A)}\theta'(g',h;\Phi')\,dg'.
\]
A configuration is \emph{admissible} if $E/F$ is unramified at all finite places, $E/\Q$ is Galois, all inert
places are odd of norm $>3$, the infinity types satisfy $m_\alpha-m_\beta>a+a'$, and $a-a'$ and $a'$ are large
compared with the real rank $p+q$ of $V'$.

\begin{definition}[Local matching; agreement of theta integrals]\label{def:families:K}
$\sfname{LocalMatching}$ asserts: for every admissible configuration there are local identifications (``real
and finite maps'') under which, at every place $v$, a common irreducible quotient of the local theta lifts of
the trivial representations of $\Uni(V_v)$ and $\Uni(V'_v)$ is matched, with matching spherical vectors and,
at the finitely many ramified split places, an explicit description by zeta integrals of Godement--Jacquet
type. $\sfname{ThetaIntegralsAgree}$ asserts: (i) for admissible configurations the theta functions and
integrals converge and are of moderate growth ($I'$ only for $\Phi'$ whose component at an auxiliary split
place is supported on matrices of full rank), and for matched data there is a unique constant $c\ne0$ with
$I(\Phi)=c\,I'(\Phi')$ for all globally matched pairs $(\Phi,\Phi')$, together with a description of the local
coinvariants; (ii)--(iii) two \emph{step relations}: removing a hyperbolic plane from $V'$ (resp.\ $V$) and
from $W$ multiplies the constant by an explicit factor built from partial Hecke $L$-values
$L^{S_0}(k,\alpha/\beta)$ at integers $k\ge2$, local epsilon factors and values of $\alpha/\beta$ on fractional
ideals, up to a non-zero constant depending only on the geometric data.
\end{definition}

\begin{theorem}[Theorem K]\Fmark\label{thm:theorem-K}
$\sfname{ThetaIntegralsAgree}$ and $\sfname{LocalMatching}$ hold.
\formal{PositivityTheory.ThKR.theoremK_head2}
\formal{UpperTheory.LiteratureTheorems.theta_integrals_agree_and_local_matching_of_literature}
\end{theorem}

\begin{proofoutline}
Two analytic inputs are proved first: the meromorphic continuation of Siegel Eisenstein series on
$\Uni(Y\oplus-Y)$ (cf.\ \cite{Langlands1976}), and the Siegel--Weil formula for a definite space $A$ of dimension
$m$ paired with $Y\oplus-Y$, $\dim Y=n$: at $s_0=(m-n)/2$ the Eisenstein series of $\Phi$ equals
$\kappa\,\mu(\mathcal D)^{-1}\int_{\mathcal D}\theta$, with $\kappa\in\{1,2\}$ \cite{Weil1965,KudlaRallis1988,Ichino2004};
the latter by Poisson summation and a comparison of Whittaker--Fourier coefficients (about $21\,000$ lines). From
fourteen local and analytic statements (local matching at split, inert and ramified places, Howe's unramified
computation \cite{Howe1979}, a modulus computation at inert places, reduction theory, existence of Weil
representations) one derives twenty-six ``leaves'' (constant terms along hyperbolic planes, local factors of the
steps, non-vanishing of the bottom constant via the doubling integral, the definite base case); the theorem then
follows by induction on $p+q$.

%P2PTR thm:theorem-K
A complete proof is given in \cref{sec:pf:thmK}.
\end{proofoutline}

\begin{remark}\label{rem:families:K}
Identities between theta integrals of complementary spaces are of the type of the first-term identities of the
regularised Siegel--Weil formula \cite{KudlaRallis1994,Ichino2004,GanQiuTakeda2014,Yamana2014}, and
$\sfname{LocalMatching}$ concerns the local theta correspondence of the trivial representation
\cite{HarrisKudlaSweet1996,KudlaSweet1997,Minguez2008}; we do not know a published statement of this form. The
second formal name takes an argument (three statements of \cref{sec:galois}) which is not used.
\end{remark}

\paragraph{Lifts and segments.}
The basic step of \eqref{eq:families:chain} is the theta correspondence between $\Uni(W)$ and a totally
definite $\Uni(V)$ with $\dim V=\dim W+t$ and $\dim W=r$.

\begin{definition}[Segments]\label{def:families:string}
For $q>0$, $z\in\C^\times$ and $n\ge1$ let $\familiesstr(q,z,n)$ be the multiset
$\{z\,q^{(n-1)/2-k}:0\le k<n\}$. A system $s$ of Hecke eigenvalues on $\Uni(V)$ (at the primes of $E$ prime to
the level) \emph{has a segment} if it admits a Satake family $B$ of rank $\dim V$ such that at all good places
$w$
\[
B(w)=B'(w)+\familiesstr\bigl(\Nm w,\chi_V(w),t\bigr),\qquad |b|\le\Nm w^{(r-1)/2}\ \text{for } b\in B'(w),
\]
where $\chi_V(w)$ is the value at $w$ of the unitary Hecke character attached to the splitting data.
\end{definition}

\begin{proposition}[Non-vanishing and injectivity of the lifts]\Fmark\label{prop:families:liftinj}
Let $D$ be a global datum and $R$ a realisation of its theta data satisfying the standing hypotheses of the
development (core conditions, good types at $g_0,g_1$, and a gap of at least $100r^4$ between consecutive
Harish-Chandra parameters). For $0\le i<m$ consider the lift from the definite $\Uni(V_i)$ to $\Uni(W_i)$. Then:
\begin{enumerate}[label=(\alph*)]
\item every non-zero Hecke eigenvector on $\Uni(V_i)$ whose system has a segment has a non-zero theta lift to
$\Uni(W_i)$;
\item a ``good'' eigenvector (one fixed by the idempotent of the type) whose system has a segment and whose lift
vanishes is zero;
\item if $x,y$ are eigenvectors at the base character $\mu_0$ with the same system, the lift of $x$ is non-zero,
$y$ is good and the lift of $y$ vanishes, then $y=0$.
\end{enumerate}
\formal{GlobalTheory.T.liftInj_and_nonVan_typeNZ_of_given}
\end{proposition}

\begin{proofoutline}
This is the principle ``pole $\Rightarrow$ non-vanishing $\Rightarrow$ injectivity'' of Rallis and
Kudla--Rallis \cite{Rallis1984inj,KudlaRallis1994,GanQiuTakeda2014,Yamana2014}, from rapid decay of cusp forms,
the doubling method, the unramified doubling computation for $\GL_n$, partial Hecke $L$-functions and Howe's
unramified statement. If the parameters have a segment, the twisted partial $L$-function factors as
$L^S_{\mathrm{tw}}(s)\prod_{i=0}^{t-1}\zeta_E^S\bigl(s+\tfrac{t-1}2-i\bigr)$; at $s_0=(1+t)/2$ the arguments
of $\zeta_E^S$ are $t,t-1,\dots,1$, so there is a simple pole, the first factor being holomorphic and non-zero
because $r+5\le t$. Through the doubling integral and a residue identity for Siegel Eisenstein series on
$\Uni(n,n)$ (about $70\,000$ lines) the pole forces a non-zero lift. The unramified correspondence,
$B(w)=B'(w)\chi_V(w)\chi_W(w)^{-1}+\familiesstr(\Nm w,\chi_V(w),t)$ when $\theta(y)\ne0$ has Satake parameters
$B'$, follows from Howe's computation \cite{Howe1979} and gives (b) and (c).

%P2PTR prop:families:liftinj
A complete proof is given in \cref{sec:pf:liftinj}.
\end{proofoutline}

\paragraph{Towers, seesaws and Bessel periods.}
The second ingredient is a separation statement for ``rows'': vectors in the multiplicity spaces of the tower
built by iterating the lifts along $(0,0)\to(1,0)\to\cdots\to(m,0)$.

\begin{proposition}[Separation of rows by theta pairings]\Fmark\label{prop:families:fs}
$\sfname{ThetaLocalInputs}$, Howe's unramified statement and the Siegel--Weil formula in the convergent range
imply $\sfname{FunctionalsSeparate}$: for every $\eps$ and every $C_F\ge6510$ there is $x_0$ such that for every
datum and size data at a scale $x\ge x_0$, every realisation at the base level and every non-zero smooth vector
$w$ of the base model, each non-zero row form built from $w$ is detected by a convergent pairing with a column
theta function: $f_{\Psi,\Phi}(\mathrm{rowForm}(w,t))\neq0$ for suitable $(\rho,\xi,\Psi,\Phi)$.
\formal{UpperTheory.Lit.functionalsSeparate_of_howeUnram_of_siegelWeil}
\end{proposition}

\begin{proofoutline}
Every node $(i,0)$, $i\le m$, carries a discrete series model; seesaw pairs give polynomial intertwiners
realising the theta pairing of vacuum vectors; the lift along each arrow converges and preserves
multiplicity-one theta vectors. For every non-zero row form some core form has non-zero period against it, and
core forms lie in the span of the reachable vectors (by Howe's unramified statement and the Siegel--Weil formula).
The decisive step is an \emph{ascent of Bessel periods}: along a column chain, non-vanishing of the bottom period
implies non-vanishing of the top Bessel pairing of some theta lift (about $34\,000$ lines); a Bessel-gap argument
concludes. The archimedean input consists of orthogonality relations and decay estimates for matrix coefficients
of discrete series of $\Uni(n+1)$- and $\Uni(n)$-type groups with interlacing parameters.

%P2PTR prop:families:fs
A complete proof is given in \cref{sec:pf:fs}.
\end{proofoutline}

\begin{remark}\label{rem:families:seesaw}
This corresponds to the seesaw step of \cite{Announcement2026}, in which the period of the iterated lift is
expressed through a Bessel period, as in the Gan--Gross--Prasad framework and the descent method
\cite{GanGrossPrasad2012,GinzburgRallisSoudry2011}; seesaw pairs and towers are classical
\cite{Kudla1994,HarrisKudlaSweet1996}, see also \cite{Li1992}. The uniqueness of Bessel models used for discrete
series with large gap is a special case of \cite{JiangSunZhu2010,SunZhu2012}, proved inside the development.
\end{remark}

\paragraph{Rationality of the theta vectors.}
For complex realisation data $X$ (a finite-dimensional space and theta maps from the rational points of the
family into it) the \emph{row test} $\sfname{RowTest}(X)$ asserts that some linear forms separate points and
take values in the base number field on all theta vectors. $\sfname{RowRat}$ asserts $\sfname{RowTest}$,
eventually in the scale and for large twisting data, at all starting configurations of the construction of
\cref{sec:families:nonempty}.

\begin{proposition}[Rationality of rows]\Fmark\label{prop:families:rowrat}
$P$, $F$, \cref{thm:theorem-K} and $\sfname{FunctionalsSeparate}$ imply $\sfname{RowRat}$.
\formal{GlobalTheory.T.test_vector_assembly_holds}
\end{proposition}

\begin{proofoutline}
By unfolding the theta integrals into orbit sums (a method going back to Rallis \cite{Rallis1984}), a flip
invariance of the twisted lifts and the assembly of their jets, the theta pairings are, up to a universal
constant $\kappa$ and roots of unity, class sums with values in the number field. Non-vanishing of Hecke
$L$-values (from the functional equation \cite{Tate1967thesis}), a Witt-tower norm formula, injectivity from
$\sfname{FunctionalsSeparate}$ and cuspidality of the lifts give the separating forms.

%P2PTR prop:families:rowrat
A complete proof is given in \cref{sec:pf:rowprops}.
\end{proofoutline}

\begin{remark}\label{rem:families:mult}
Several rows also use an equality of cuspidal multiplicities for $\Uni(Y_0)$ and $\Uni(W_i)$ when $Y_0$ and $W_i$
are isometric at all finite places, proved by a simple trace formula with pseudo-coefficients at infinity and
types at a finite place \cite{DeligneKazhdanVigneras1984,Arthur1989,ClozelDelorme1984}; it is a comparison of
multiplicities, not multiplicity one in the sense of \cite{AizenbudGourevitchRallisSchiffmann2010}.
\end{remark}

% ---------------------------------------------------------------------------
\subsection{Rows and linked bundles}\label{sec:families:rows}

\paragraph{Families and closedness.}
We recall the notions of \cref{sec:top,sec:stepanov} in the form used here.

\begin{definition}[Families; closed properties]\label{def:families:closed}
A \emph{family} $\mathcal F\in\sfname{FamP}(\eps)$ consists of a threshold $x_1>1$, index sets $I_x$
($x\in\R$), members $(D_{x,i},L_{x,i})$ whose size data have scale $(x,\eps)$ for $x>x_1$, extra data
$\mathrm{ex}(x,i)$ (levels, theta data, twists and a \emph{norm choice} $\mathrm{nch}_{x,i}$), and lists of
family constants; a \emph{family class} is a predicate on families. For a class $\mathcal C$ and a property
$P(\mathcal F,C_H)$, $\familiesClosed(\mathcal C,P)$ means: there is $\eps_1>0$ such that for every
$\eps\in(0,\eps_1]$ there is a staged threshold $K$ with $P(\mathcal F,C_H)$ for all
$\mathcal F\in\mathcal C\cap\sfname{FamP}(\eps)$ and all $C_H$ for which the constants of $\mathcal F$ are large
with respect to $(K,C_H)$. For a class $\mathcal A$, $\familiesDiii=D3(3,\mathcal A)\subseteq\mathcal A$ is the
subclass of \emph{designed} families (twist goodness $\sfname{TwGood}$, smoothness, divisibility, tameness and
box conditions, the size conditions $\sfname{TuSize}$, $\sfname{DGamSize}$ and a choice at level $3$).
\end{definition}

A closed property says nothing about classes without large members; existence of large members is a separate
statement (the field $\mathsf{hEx}$ of the package, \cref{sec:families:nonempty}).

\paragraph{Points, bundles and rows.}
The members of the families in the package are produced by a construction, and their properties are proved
uniformly for all its outputs. A \emph{context} fixes a scale and constants (``eventually in the scale'' means:
for small $\eps$, staged large constants and then large $x$). A \emph{base} fixes ranks, host, hermitian spaces,
the words $u$ and a coefficient field; a \emph{frame} fixes Hecke characters of $E$; twisting data are described
in \cref{def:families:twist}. A \emph{core} is the datum at one level of a construction, and a \emph{slot} is a
predicate on cores. A \emph{bundle} over a datum consists of a reciprocity record, spectral ties, a row-side
family, square-integrable columns and an \emph{indefinite realisation} (the theta package on the indefinite side
with its lattices and metrics); a \emph{linkage} relates bundles to points; and a \emph{selection} $\sigma$ for a
class $\mathcal A$ chooses a bundle $\sigma(\mathcal F,x,i)$ over each member $D_{x,i}$, $\mathcal F\in\mathcal A$.

\begin{definition}[Rows; linked statements]\label{def:families:rows}
Let $Q$ be a slot, $\mathrm{Lk}$ a linkage and $H(\mathcal A,\sigma)$ a property of classes with selections.
$\familiesRow(Q,\mathrm{Lk},H)$ means: there is a property $E$ of bundles which holds, eventually in the scale,
for all bundles linked to points of $Q$, such that $H(\mathcal A,\sigma)$ holds for every class $\mathcal A$
and every selection $\sigma$ whose bundles are built from points of $Q$ and satisfy $E$. A \emph{parameter}
$p$ consists of a norm-choice extension (extra data demanded at each point), a spectral extension and a
geometric predicate; it determines slots $\sfname{IsGoodE}(p)\subseteq\sfname{IsBuiltE}(p)$ of good and built
points and a linkage $\sfname{BundleLinkage}(p)$. For a property $r$ of slots and linkages,
$\familiesLinked(p,r)$ means: there is a slot $R$ which, eventually in the scale and for large twisting data,
contains every built point, such that $r\bigl(\sfname{IsGoodE}(p)\wedge R,\ \sfname{BundleLinkage}(p)\bigr)$.
\end{definition}

Linked rows are closed under conjunction, and passing to a stronger norm-choice extension (fewer built and good
points) preserves them. The structure $\sfname{AllRowsAlt}$ is the conjunction of twenty-six linked rows which
correspond to the fields of $\sfname{HypsAlt}$ concerning the selection (\cref{tab:families:fields}).

\begin{proposition}[The row package]\Fmark\label{prop:families:allrows}
Assume global Artin reciprocity, the Siegel--Weil formula in the convergent range (\cref{lit:families:sw}), five
classical statements on real unitary groups and their automorphic forms, $\sfname{SurjectiveOntoX}$ and
$\sfname{LambdaLowerBound}_2$ (\cref{sec:families:nonempty}), and $\sfname{RecTExists}$. Let $p$ be a parameter
attached to a norm-choice extension and $\mathcal N$ a predicate on classes. If the twelve open parameter
statements hold for $p$ (\cref{thm:families:open}) and $p$ is linked for the covolume row, the functional row and
the norm row of $\mathcal N$, then $\sfname{AllRowsAlt}$ holds for $p$, with norm predicate $\mathcal N$ and
definite-side condition $\sfname{DefiniteSideCondition}$.
\formal{GlobalTheory.Fc.AllAtMemberStage0.allRows_of_open0_alt}
\end{proposition}

\begin{proofoutline}
Six rows are hypotheses, and several hold for every slot and linkage. The substance lies in the local rows from
the Siegel--Weil formula (\cref{thm:theorem-S}), the service row (CM theory, below), the counting row
(\cref{thm:counting}), the period rows (\cref{prop:families:periods}) and the definite rows.

%P2PTR prop:families:allrows
A complete proof is given in \cref{sec:pf:rows}.
\end{proofoutline}

\paragraph{The Siegel--Weil formula and its consequences.}
The archimedean size of the construction is controlled through the following classical theorem, which is a
field of $P$ and is proved in the development.

\begin{littheorem}[Siegel--Weil formula, convergent range {\cite{Siegel1935,Weil1965}}]\label{lit:families:sw}
Let $E/F$ be a CM extension, $J$ a non-degenerate hermitian $M\times M$ matrix over $E$ and $H=\Uni(J)$. Then
$H(F)\backslash H(\A_F)$ has finite volume, and if $2n<M$, $T\in M_n(E)$ is hermitian with $\det T\ne0$, and
$\Phi=\Phi_\infty\otimes\bigotimes_v\Phi_v$ on $(\A_E^M)^n$ has $\Phi_\infty$ smooth and compactly supported and
$\Phi_v$ Schwartz--Bruhat (almost all standard), then the mean over $H(F)\backslash H(\A_F)$ of
$h\mapsto\sum_{X\in(E^M)^n,\ (h_J(x_a,x_b))=T}\Phi(h^{-1}X)$ equals $g_\infty(T)\prod_{v<\infty}g_v(T)$, where
$g_\infty,g_v$ are the densities of the push-forwards of $\Phi_\infty,\Phi_v$ under the Gram map (for self-dual
measures).
\end{littheorem}

The formal statement is the special case of Weil's theorem with compactly supported $\Phi_\infty$; regularised
versions \cite{KudlaRallis1988,KudlaRallis1994,Ichino2004,Yamana2014} are not needed.

\begin{proposition}[Consequences of the Siegel--Weil formula]\Fmark\label{prop:families:volratio}
\Cref{lit:families:sw} implies: (i) a mass-ratio identity for the hermitian volumes of every datum;
(ii) the Siegel--Weil volume identities $\sfname{SWXS}$ for the CM and hermitian steps of every Hida setting;
(iii) a bound $\varepsilon_{\mathrm{NG}}\,d_ud_{u'}/D_*$ for the non-generic terms of the orbital expansions
below ($\sfname{RowNGP}$); and (iv) the volume recursion $\sfname{VolumeStep}$: an automorphic instance $I$ of a
unitary group of rank $n\ge3$ and the instance $I'$ of rank $n-1$ given by the stabiliser of a vector $x$ with
$\langle x,x\rangle=\pm1$ satisfy $\vol(I)/v_U(n)^{f}=\mathrm{massRatio}(E,n)\,\vol(I')/v_U(n-1)^{f}$, where
$f=[F:\Q]$, $v_U$ is explicit and $\mathrm{massRatio}(E,n)$ is the right-hand side in \cref{prop:families:mass}.
\formal{GlobalTheory.GaussMean.volumeStep_of_siegelWeil}
\end{proposition}

\begin{proofoutline}
For $n=1$ the formula gives Siegel's mean-value formula for one vector; a limiting argument extends it to
general test functions, and unfolding the mean of a Gaussian over the orbit of a unit vector gives (iv). The
bound (iii) also uses the uniqueness of Bessel models for discrete series with large gap.

%P2PTR prop:families:volratio
A complete proof is given in \cref{sec:pf:rows}.
\end{proofoutline}

\paragraph{Approximate orthogonality.}
For a member of a family and a word $u$, the spectral data of the selection provide a Hida setting (an
automorphic instance of a unitary group over $E$ with a level and discrete series data at infinity) and a space
$C_u$ of automorphic forms of dimension $d_u$. Let $D_*$ be the \emph{main term} (formal degree times volume times
the value at $1$ of the level idempotent). For words $u,u'$ the operator relating their flag data has an orbital
expansion over pairs of row lattices $\Lambda,\Lambda'$ and classes of rational isometries $\gamma$ with
$\gamma\Lambda'=\Lambda$; let $K_0(u,u')$ be its leading coefficient (identity, generic and non-generic terms).

\begin{theorem}[Theorem S]\Fmark\label{thm:theorem-S}
Assume $\sfname{SWXS}$, $\sfname{RowNGP}$ and finite-dimensionality of the row spaces. Let $\mathcal C$ be a
family class with a spectral selection such that all selected data are locally tied, the rank row
$\familiesClosed(\mathcal C,\sfname{RanksF})$ holds, and the loop rows hold. Then, closedly on $\mathcal C$,
eventually in $x$: for all $i\in I_x$ and all distinct words $u\ne u'$ of $D_{x,i}$ (under the admissible
bounds of the selection),
\[
\|K_0(u,u')\|^2\ \le\ \bigl(1+b^{-6}\bigr)\,\frac{d_u\,d_{u'}}{D_*},\qquad b=e^{x^\eps}.
\]
\formal{GlobalTheory.CarrierGap.offDiag}
\end{theorem}

\begin{proofoutline}
The proof (about $193\,000$ lines) first establishes the bound pointwise, for every datum and pair of distinct
words satisfying size, count, gap and local-tie hypotheses, the trace frame and the loop-row inequalities; the
same bound holds for the Hilbert--Schmidt mass of the off-diagonal block. Its nineteen hypotheses come from local
models, density inputs, $\sfname{RowNGP}$, $\sfname{SWXS}$, convergence and folding of the orbital integrals,
Wu decay and a counting formula. The key analytic input is a lower bound for ``Poincar\'e amplitudes'' (about
$67\,000$ lines), deduced from a compact-model theorem, a model lemma for oscillatory orbital integrals (stationary
and non-stationary phase) and the uniqueness of Bessel models in large gap. A transport lemma passes to families.

%P2PTR thm:theorem-S
A complete proof is given in \cref{sec:pf:thmS}.
\end{proofoutline}

\begin{remark}\label{rem:families:random}
For orthogonal projections $P,P'$ onto a fixed subspace of dimension $d$ and a uniformly random subspace of
dimension $d'$ of $\C^{D}$ one has $\mathbb E\,\|PP'\|_{\mathrm{HS}}^2=dd'/D$. Theorem~S and $\sfname{LinkS}$
(\cref{cor:families:linkS}) say that the spaces attached to distinct words are, up to a factor $1+o(1)$, as
orthogonal as independent random subspaces of dimension $D_*$: the approximate orthogonality of the columns of
the integral structure in \cite{Announcement2026}, described there through a relative trace formula and projected
Poincar\'e series. We do not know of a published result of this form.
\end{remark}

\paragraph{Counting.}
Let $I$ be an automorphic instance of a unitary group of rank $n$ over a CM extension $E/F$, $f=[F:\Q]$, with
root discriminant $r_d$, discrete series data with Harish-Chandra parameters $\lambda_v$ at the infinite places,
a level and a type at an auxiliary prime $g_0$. Let $D_{\mathrm{count}}$ be the dimension of the space of
automorphic forms of the given type, and $D_*$ the main term.

\begin{theorem}[Automorphic counting with exponentially small error]\Fmark\label{thm:counting}
Assume the volume recursion $\sfname{VolumeStep}$ of \cref{prop:families:volratio}. For every $C_g$ there is
$n_0$ such that the following holds for every instance as above of signature $(p,q)$ and every $b\ge2$ with
$n\ge n_0$, $q^2\le n$, $\log r_d\le e^n$, $\log f\le\sqrt n$, $\log b\le\exp(\tfrac1{16}nf\log n)$,
$|\lambda_v(a)|\le b^2$, a level of size $1\le g\le b^{C_g}$ and a type prime $g_0>4n$ split completely in
$E$: if all $\lambda_v$ have gaps at least $8\cdot10^{14}(C_g+4)\,n^{19}f^5\log(b\,r_d)$, then the space of
automorphic forms is finite-dimensional and
\[
\bigl|D_{\mathrm{count}}-D_*\bigr|\ \le\ e^{-\frac1{16}nf\log n}\,D_* .
\]
\formal{GlobalTheory.CountFinal.countingTheorem_of_volumeStep}
\end{theorem}

\begin{proofoutline}
A pre-trace formula argument (about $36\,000$ lines). A compression formula for the lowest $K$-type gives
decay and a pointwise bound for discrete series matrix coefficients with large gap. By induction on $n$,
$\sfname{VolumeStep}$ gives an absolute volume bound and bounds for the probability of isotropic and of non-zero
flags. These control the contributions of the near-identity, far, rare-flag and deep regions of the geometric
side, so that $D_{\mathrm{count}}$ equals the identity contribution $D_*$ up to the stated error.

%P2PTR thm:counting
A complete proof is given in \cref{sec:pf:rows}.
\end{proofoutline}

\begin{remark}\label{rem:families:counting}
Limit multiplicity and Plancherel equidistribution for families of automorphic representations are known with
power savings in the level \cite{Shin2012,ShinTemplier2016}; \cref{thm:counting} is of a different nature, with a
relative error exponentially small in $nf\log n$ under a gap condition. Applied to the label data at the type
primes it yields the counting row of the package, with relative error $\exp(-c\,rf\log r)$.
\end{remark}

\paragraph{Masses.}
The mass of a hermitian lattice is the sum of $1/|\Aut|$ over a set of representatives of the classes in its
genus.

\begin{proposition}[Siegel mass ratio]\Fmark\label{prop:families:mass}
Let $K$ be a CM field such that $K/K^+$ is unramified at all finite primes, $d=[K^+:\Q]$,
$v=|\mathrm{disc}\,K|^{1/(2d)}/2$ and $M\ge3$. Let $\mathrm{mass}_M$ be the mass of the standard lattice
$\OO_K^M$ with the form $\sum_i\bar x_iy_i$. Then
\begin{gather*}
\frac{\mathrm{mass}_M}{\mathrm{mass}_{M-1}}=\frac{\bigl((M-1)!\,(v/\pi)^M\bigr)^{d}}{(2v)^{d/2}}
\prod_{\mathfrak P}e_M(\mathfrak P),\\
e_M(\mathfrak P)=\bigl(1-(-1)^M\Nm\mathfrak P^{-M/2}\bigr)^{-1}\ \text{if }\bar{\mathfrak P}=\mathfrak P,\qquad
e_M(\mathfrak P)=\bigl(1-\Nm\mathfrak P^{-M}\bigr)^{-1/2}\ \text{if }\bar{\mathfrak P}\ne\mathfrak P,
\end{gather*}
the product running over the primes of $K$; it equals $L(M,\eta^M)$ for the quadratic character $\eta$ of
$K/K^+$.
\formal{GlobalTheory.Mass.massRatio_all_unramified_cmFields}
\end{proposition}
%P2PTR prop:families:mass
\begin{proof}See \cref{sec:pf:bundles}.
\end{proof}


This is a mass-ratio form of Siegel's mass formula for unimodular hermitian lattices
\cite{Siegel1935,Prasad1989,GanHankeYu2001}. It enters the volume recursion, the class-number and genus-mass
rows, the explicit normalising constant of the row counts (a unitary analogue of Siegel's mass formula for the
row lattices, proved from $\sfname{SWXS}$), and the family statements of \cref{prop:families:famgl}.

\paragraph{Periods, levels and CM theta functions.}
The remaining rows concern the arithmetic of the period entries. At good points the data are assembled from
``true columns'', and six rows hold for every slot and linkage.

\begin{proposition}[Period entries and level valuations]\Fmark\label{prop:families:periods}
At bundle points over true columns, eventually: the reference-chain evaluation maps are injective with integral
period points; the periods at $\mu_0$ and at twisted characters are values of elements $y\in\mathit{LL}$ with
$y=0$ or $v_w(y)$ bounded below explicitly at every place $w$ (outside the twisting primes in the twisted case);
the real-place ties hold; and at every place $w$
above the level prime $T_u$ of a word $u$, if $c\in\mathit{LL}$ has $v_w(c)\ge e_w\,m\,E_T(u,w)$, then
$c\cdot\mathrm{col}_T(u,\varphi)$ lies in the lattice $S_m(w)$ for all columns $\varphi$ in the relevant range.
\formal{GlobalTheory.T.period_entries_tie_and_level_valuations_at_built}
\end{proposition}
%P2PTR prop:families:periods
\begin{proof}See \cref{sec:pf:rows}.
\end{proof}


The last statement (about $35\,000$ lines) is a divisibility at the level primes whose order is proportional to
the chain length $m$. The \emph{service row} concerns the theta kernels of the construction, whose theta
component is built from CM elliptic curves: the lattice $\delta\cdot\mathfrak c\subset\C$ of a non-zero ideal
$\mathfrak c$ of $\OO_K$ of norm prime to $2d$, where $\delta^{12}$ is the discriminant of the lattice
$\mathfrak c$, with a marked $2$-torsion point. From global Artin reciprocity and the Chebotarev density theorem
the development proves that these curves have good models over $\Q_\ell^{\mathrm{ur}}$ for $\ell\nmid6d$
\cite{Deuring1941,SerreTate1968,Shimura1971}, that at primes $\ell\ge5$ inert in $K$ their formal groups become
Lubin--Tate after a unit twist \cite{Honda1970}, and that their prime-to-$\ell$ torsion points are integral; it
deduces the algebraicity of the Taylor coefficients of the normalised theta functions at torsion points and the
integrality of the kernels at inert primes.

\begin{proposition}[Chowla--Selberg]\Fmark\label{prop:families:cs}
Let $K$ be an imaginary quadratic field of discriminant $-d$ with $d>4$, $\chi$ its Kronecker character and
$h$ its class number. Then $L(1,\chi)=\pi h/\sqrt d$, and for points $\tau_{\mathfrak c}$ of the upper half plane
representing the ideal classes $\mathfrak c$ of $K$,
\[
\frac1h\sum_{\mathfrak c}\log\bigl(4\pi\,\Im(\tau_{\mathfrak c})\,|\eta(\tau_{\mathfrak c})|^4\bigr)
=\gamma+\log\frac{2\pi}{\sqrt d}-\frac{L'(1,\chi)}{L(1,\chi)} .
\]
\formal{GlobalTheory.chowla_selberg}
\end{proposition}

This is the formula of Chowla and Selberg \cite{ChowlaSelberg1967}. The development uses it to evaluate a
normalised mean over the conjugates of the logarithmic periods of the lattices $\delta\cdot\mathfrak c$, equal to
$\gamma+\log(2\pi/r_E)-L'(1,\chi)/L(1,\chi)$ ($r_E$ the root discriminant of $E$), which enters the kernel bounds.

% ---------------------------------------------------------------------------
\subsection{The Galois link}\label{sec:families:galois}

The norm row $\mathsf{nRest}$ of the package is the closedness of $\sfname{LinkN}$ on $\familiesDiii$. It bounds
the failure of divisibility at the twisted places by Galois-cohomological quantities.

\paragraph{The Galois side of a datum.}
The Galois data of $D$ consist of an algebraic closure $\Omega$ of $E$, for each prime $\ell$ the set $S_\ell$ and
the cyclotomic character of $G_{E,S_\ell}$, and a finite set of \emph{local indices} $\iota$. A local index carries
a prime $\ell=\ell_\iota$, a lift index $1\le j_\iota\le m$, a twisted place $w_\iota$ of $\mathit{LL}$ above $\ell$,
a complete discrete valuation ring $\OO_\iota$ finite over $\Zl$, a finite set $X_\iota$ of eigensystems with
subsets $Z^{\mathrm{good}}_\iota\subseteq Z_\iota$, a local Hecke algebra $T_\iota\subseteq\OO_\iota^{X_\iota}$ acting
faithfully on a free $\OO_\iota$-module $M_\iota$, Galois characters, and finite sets $S_\iota,Q_\iota$ of auxiliary
split places. Lengths are $\OO_\iota$-lengths divided by the length of $\OO_\iota/\ell$ (an infinite length counts as
$0$). Then $\eta_{\mathrm{loc}}(\iota)$ is the length of a quotient $M^{\mathrm{up}}_\iota/M^{\mathrm{low}}_\iota$ of
two lattices, and $\eta_j(w)=\sum\eta_{\mathrm{loc}}(\iota)$ over the $\iota$ with $(j_\iota,w_\iota)=(j,w)$ is the
divisibility defect of the $j$-th lift at $w$; $v_{cc}(\iota,\pi)$ is the length of $\OO_\iota/\mathfrak c_\pi$, with
$\mathfrak c_\pi$ the image at $\pi\in X_\iota$ of the congruence ideal of $T_\iota$ relative to $Z_\iota$, and
$\mathrm{mult}(\iota,\pi)$ the multiplicity of $\pi$ in $M_\iota$; for Galois representations
$\rho_{\iota,\pi}:G_E\to\GL_r(\OO_\iota)$, $s_\Pi(\iota,\rho,\pi)$ is the length of a Selmer group with Bloch--Kato
conditions at $\ell$ \cite{BlochKato1990}, relaxed at $S_\iota\cup Q_\iota$; and $\mathrm{sh}(\ell,e)$ is the
$\Zl$-length $\mathrm{sh}_\ell(E,e)$ of $\Sha^1(G_{E,S_\ell},\Ql/\Zl(-e))$ (\cref{sec:prelim}) when it is finite, and
$VV(\ell)=\sum_{0<e<t}\mathrm{sh}(\ell,e)$.

\begin{definition}[The link $\sfname{LinkN}$]\label{def:families:linkN}
For size data $L$ of $D$ and $1\le j\le m$ put
\[
\eta\Sigma(L,j)=\sum_{\substack{w\in\supp(L)\\ \mu_w\ne\mu_0}}\frac{\log\Nm w}{[\mathit{LL}:\Q]}\;e(w)\,\eta_j(w),
\]
with $e(w)$ the ramification index. A family $\mathcal F$ satisfies $\sfname{LinkN}$ if there are $C,x_0$ such
that for all $x\ge x_0$ and $i\in I_x$ for which all the groups $\Sha^1(G_{E,S_\ell},\Ql/\Zl(-e))$, $0<e<t$, are
finite, and all $1\le j\le m$: if $\sum_\ell VV(\ell)\log\ell$ converges, then
\begin{equation}\label{eq:families:linkN}
\eta\Sigma(L,j)\ \le\ d_{Z_0}(j)\Bigl(C\bigl(f\,m\,b\,\Lambda^2+(\log x)^2N_c\bigr)
+\sum_{\ell}\log\ell\sum_{0<e<t}\mathrm{sh}(\ell,e)\Bigr).
\end{equation}
$\sfname{NL}(\mathcal C)=\familiesClosed(\mathcal C,\sfname{LinkN})$.
\end{definition}

\paragraph{Inputs.}
The Galois-theoretic inputs are fields of $P$, proved in the development: local finiteness, Tate duality and
the Euler characteristic formula, for finite extensions of $\Ql$ and for decomposition groups
\cite{Tate1963,Serre1997,NeukirchSchmidtWingberg2008}; a degree-one Chebotarev theorem \cite{Chebotarev1926}; and
the package $\sfname{Hypotheses}_2$ (Fontaine's $B_{\mathrm{cris}}$ \cite{Fontaine1982}, Artin reciprocity, the
Bloch--Kato conditions \cite{BlochKato1990}, local Galois cohomology, and the integral comparison of Bhatt--Scholze
\cite{BhattScholze2022,BhattScholze2023}). A \emph{norm choice} attaches to each twisted place $w$ a twisting
character $\psi'_w$ of a ray class group, with values in a finite extension of $\Zl$. A family is
$\sfname{TwGood}$ if eventually its norm choices are conjugation-symmetric and all $\psi'_w$ are \emph{good}
(they avoid an explicit finite set of places and satisfy index bounds, lattice inclusions and three Bloch--Kato
reduction conditions). The \emph{Galois preconditions} of a datum with norm choice are eighteen conditions, among
them admissibility at all local indices, irreducibility of the residual representations of rank $r$,
unramifiedness of the characters, a depth-zero condition and uniqueness of localisations; $\sfname{Tie4}$ ties the
characters at the type primes to de Rham local Galois representations of prescribed Weil--Deligne type.

\begin{theorem}[The Galois link]\Fmark\label{thm:linkN}
\begin{enumerate}[label=(\alph*)]
\item Assume $\sfname{Hypotheses}_2$ and Chebotarev's theorem and local Galois cohomology for decomposition groups.
Let $\mathcal F\in\sfname{FamP}(\eps)$ be $\sfname{TwGood}$, and suppose that for all $x>x_1$ and $i\in I_x$ the
datum $D_{x,i}$ satisfies the Galois preconditions for its norm choice and $\sfname{Tie4}$. Then
$\mathcal F$ satisfies $\sfname{LinkN}$.
\item Assume in addition the remaining Galois-theoretic statements derived from $P$ (the third-layer inputs,
compatibility of the Galois representations attached to polarized cuspidal representations, rationality and
integrality of twist series, depth-zero support). For
every $h_L:\sfname{RecTExists}$ and $c\in\R$ the parameter $p=\sfname{parametersOfDependence}(h_L,c)$ of
\cref{sec:families:nonempty} satisfies
\[
\begin{aligned}
&\familiesLinked\bigl(p,\ \familiesRow(\cdot,\cdot,\ \mathcal A\mapsto\familiesClosed(D3(3,\mathcal A),
\sfname{BoundAtTwoThree}))\bigr)\\
&\qquad\Longrightarrow\
\familiesLinked\bigl(p,\ \familiesRow(\cdot,\cdot,\ \mathcal A\mapsto\sfname{NL}(D3(3,\mathcal A)))\bigr).
\end{aligned}
\]
\end{enumerate}
\formal{UpperTheory.GalC.linkN_of_twGood_galoisPreconditions}
\formal{UpperTheory.GalC.final_alt_alt}
\end{theorem}

\begin{proofoutline}
(a) The proof (about $221\,000$ lines) runs as follows.
\begin{enumerate}[label=\arabic*.]
\item \emph{Cuts.} Each member is replaced by a \emph{good cut}: keep the local indices with
$Z^{\mathrm{good}}_\iota\ne\emptyset$ and restrict to eigensystems on which an endomorphism of $M_\iota$ commuting
with $T_\iota$ (from central idempotents of a group algebra) acts by $\#G$ on the good eigenspaces and by $0$
outside. The existence of good cuts with the ``norm-choice facts'' (a hypothesis $\sfname{CutGivesAlt}$ of the
first formal statement) is proved from Chebotarev's theorem and a vanishing statement: at an auxiliary place
$v\in S_\iota$ at which the residual representation at a tame generator has no eigenvalue equal to the residual
character value, the local $H^1$ and $H^2$ of all Tate twists vanish, by local duality and the Euler
characteristic formula (as for Taylor--Wiles primes \cite{TaylorWiles1995,DarmonDiamondTaylor1997}).
\item \emph{Local bounds} (about $163\,000$ lines, after uniform constants are fixed). For each twisted local index
and good eigensystem
\begin{equation}\label{eq:families:bd}
v_{cc}(\iota,\pi)\ \le\ s_\Pi(\iota,\rho,\pi)+\beta(\iota)+VV(\ell)+v_\ell(h_E)+\ell^{-2},
\end{equation}
with $\beta=0$ for $\ell$ beyond a threshold outside a small exceptional set (a congruence bound from local
Bloch--Kato theory for twisted lattice cohomology); $\beta\le Cfnt^2n_V\Lambda^2$ plus local lengths in the middle
range (lifting theorems for crystalline lattices with Hodge--Tate weights in $[-h,h]$, $h<(\ell-1)M$, through
Wach-type $(\varphi,\Gamma)$-modules and Herr cocycles \cite{Wach1996,Berger2004,Herr1998,FontaineLaffaille1982});
explicit bounds for $\ell\in\{q,d\}$, the small primes and $\ell=2,3$ (\cref{lem:families:at23}); and a bound by a
paired character condition at the type primes. Moreover $\eta_{\mathrm{loc}}(\iota)\le\sum_{\pi}\mathrm{mult}(\iota,\pi)
\bigl(v_{cc}(\iota,\pi)+v_\ell(\#G)\bigr)$, the multiplicities summing to at most $d_{Z_0}(j)$.
\item \emph{Summation.} With $\sum_{w\mid\ell}e(w)\log\Nm w/[\mathit{LL}:\Q]\le\log\ell$ and the killing condition
(finiteness and smallness of the Selmer groups of the good eigensystems) one gets \eqref{eq:families:linkN} for the
cut family with constant $2000(|C|+1)^3$; the $\eta$-invariants of the original members exceed those of the cut
members by a controlled amount, which gives $\sfname{LinkN}$ for $\mathcal F$.
\end{enumerate}
(b) The Galois preconditions and $\sfname{Tie4}$ are proved at all members of the construction (from the
third-layer inputs, base change, depth-zero support \cite{MoyPrasad1994} and Green's labelling of characters),
and transported to all bundles linked to good points; then (a) applies to every family in $D3(3,\mathcal A)$ for
every class $\mathcal A$ selected from good points.

%P2PTR thm:linkN
A complete proof is given in \cref{sec:pf:linkN}.
\end{proofoutline}

\begin{lemma}[The bound at $2$ and $3$]\Fmark\label{lem:families:at23}
Assume $\sfname{Hypotheses}_2$. Then every family satisfies $\sfname{BoundAtTwoThree}$: eventually, for every local
index $\iota$ with $\ell_\iota=2$, or $\ell_\iota=3$ and $3$ not a small prime, every good $\pi$ and every
admissible setup (local Hodge--Tate, Fontaine and comparison data, and hub data with killing conditions),
\eqref{eq:families:bd} holds with $\beta=Cfmb\log x+Cx^C\Lambda^2$. Consequently every parameter $p$ is linked
for $\familiesRow(\cdot,\cdot,\mathcal A\mapsto\familiesClosed(D3(3,\mathcal A),\sfname{BoundAtTwoThree}))$.
\formal{UpperTheory.N.at23_family_alt}
\end{lemma}

\begin{proof}[Proof sketch]
For $\ell=2$, $\sfname{Hypotheses}_2$ gives a local bound with explicit constants ($6000$ and $30000$); for
$\ell=3$ a similar bound holds, using that $3\ne q,d$ for every host. The maximum of the two constants and of the
two thresholds works for both primes. A property of all families is closed on every class; a row whose target
holds for all selections holds with the trivial bundle property; and a row for the good slot is linked with the
trivial auxiliary slot.
\end{proof}

\begin{remark}\label{rem:families:linkN}
(1) The local bound \eqref{eq:families:bd} bounds the length of a congruence module by that of a Selmer group,
which is the shape of the numerical criterion of Wiles and Lenstra \cite{Wiles1995,Lenstra1995} (an interpretation;
the development does not refer to it). (2) The twisting characters of the norm choice make the Selmer groups of the
good eigensystems finite and small (``killing by twisting'' in \cite{Announcement2026}). What remains is
$\sum_\ell\log\ell\sum_{0<e<t}\mathrm{sh}_\ell(E,e)$, which involves only the host field; it is bounded by the
window bound (\cref{sec:window}) through the hypothesis $\sfname{HypV}$ of $\sfname{TopG}$, which also supplies the
finiteness and summability assumed in \cref{def:families:linkN} (\cref{sec:stepanov}). (3) By
\cref{lem:families:at23} the hypothesis of \cref{thm:linkN}(b) holds for every parameter; in the proof of (b) it
is used only through its auxiliary slot, the bound at $2$ and $3$ needed in step~2 of (a) being re-derived from
$\sfname{Hypotheses}_2$.
\end{remark}

% ---------------------------------------------------------------------------
\subsection{The remaining links}\label{sec:families:links}

The links other than $\sfname{LinkN}$ are proved in the proof of \cref{thm:shape} from the fields of the package.
We describe the four main ones and the properties feeding the global theorem.

\paragraph{Irreducibility from Green's theorem.}
For a type prime $g=g_k$ ($k=0,1$) and $N=r$, let $e_U\in\C[\GL_r(\F_g)]$ be the Gelfand--Graev idempotent
attached to a non-trivial character $\psi$ of $\F_g$, and $e_J$ the analogous idempotent of a ``Jacquet-type''
subgroup of the unipotent radical determined by $(k,m)$. A \emph{label} is a monic $f\in\F_g[X]$ of degree $r$
with $f(0)\ne0$; the type data of $D$ attach to it an element $e_{\mathrm{Ser}}(f)$ of the group algebra (playing
the role of the central idempotent of the Harish-Chandra series labelled by $f$), and one puts
$\mathrm{wt}(a,f)=\Re\bigl(a\,e_{\mathrm{Ser}}(f)\bigr)(1)/\Re\,a(1)$.

\begin{proposition}[Label weights]\Fmark\label{prop:families:admLq}
Assume Green's theorem (\cref{def:green}). Then for every datum $D$ with normalized type data, $k\in\{0,1\}$ and
$g=g_k$: $\mathrm{wt}(e_U,f)\ge0$ and $\mathrm{wt}(e_J,f)\ge0$ for all labels;
$\mathrm{wt}(e_U,f)=\prod_{P\mid f}(g^{\deg P}-1)^{-1}$ (product over the monic irreducible factors) for
squarefree $f$; $(1-g^{-1})\,\mathrm{wt}(e_U,f)\le\mathrm{wt}(e_J,f)$; and
$\mathrm{wt}(e_J,f)\le(1+g^{-1})^r\,\mathrm{wt}(e_U,f)$ for squarefree $f$.
\formal{GlobalTheory.Top.Lq.admissibleLabelWeights_of_greenCuspidalParametrisation}
\end{proposition}
%P2PTR prop:families:admLq
\begin{proof}See \cref{sec:pf:top-shape}.
\end{proof}


For squarefree $f=P_1\cdots P_s$ the series of $f$ contains exactly one generic representation; it occurs once in
the Gelfand--Graev representation \cite{GelfandGraev1962,Zelevinsky1981}, and by Green's theorem its degree is
$\prod_{i\le r}(g^i-1)/\prod_j(g^{\deg P_j}-1)$, which explains the formula.

For a member $D$ let $D_{\mathrm{all}}$ be the number of labelled indices (label tuples at $g_0,g_1$) of its rows
data and $D_{\mathrm{rows}}=(1-\delta_p)D_{\mathrm{all}}$ the number of those coming from rows. $\sfname{LinkIRR}$
asserts that eventually $0\le\delta_p(D_{x,i})\le\exp(-x^{\eps_z/7})$, $\eps_z=\eps/J$.

\begin{theorem}[Irreducibility link]\Fmark\label{thm:linkIRR}
Assume Green's theorem. Let $\mathcal A$ be a class with a twist selection and $c_{Ic}>0$, and suppose that,
closedly on $D3(3,\mathcal A)$, eventually every member admits an all-rows record $A$ (labels at the type primes,
an injection of the rows into the labelled indices, and idempotents) with injective good labels, the join
property, and counts of all cylinders within relative error $\exp(-c_{Ic}\,rf\log r)$ of the expected values.
Then $\sfname{LinkIRR}$ is closed on $D3(3,\mathcal A)$.
\formal{GlobalTheory.Top.Irr.closed_linkIRR_of_join_count}
\end{theorem}

\begin{proofoutline}
The tied records of the selection provide normalized type data, so \cref{prop:families:admLq} gives the
weight properties at every selected member. With the join and a range condition they give non-negative weights
and a bound on the share of label tuples that are not good; combined with the counts and the size facts this
yields $0\le\delta_p\le\exp(-x^{\eps_z/7})$ once $x^{\eps_z}\ge3/2$. The count hypothesis is the counting row of
the package, obtained from \cref{thm:counting} (\cref{rem:families:counting}).

%P2PTR thm:linkIRR
A complete proof is given in \cref{sec:pf:top-shape}.
\end{proofoutline}

In the proof of \cref{thm:shape} Green's theorem is used only here.

\paragraph{The local lattice link.}
$\sfname{LinkG}$ asserts that there are $C$ and a function $\delta(x)\to0$ such that eventually, for all members
$(D,L)$: (i) at every place $w$ of $\mathit{LL}$ not above the level prime $T_u$ of a word $u$, the lattice $N_u$
is locally contained in an explicit lattice and its local length satisfies
$\mathrm{len}_{MN}(u,w)\le d_u(\tau(w)+\iota(w))$ for local weights $\tau,\iota$ of $L$ (vanishing where no
served orbit exists); (ii) the served orbits $t_0$ at $p$ satisfy $S_v\le2m\,N_v\,a/p$; (iii) at every place
$w$ above a prime $p\le X''(n)$ in the inert classes of the host, $0\le\mathrm{gain}(j,w)\le4|TT|\,a/p$ and
\begin{equation}\label{eq:families:linkG}
\sum_{j=1}^{m}\sum_{t\in TT_w}\bigl(\Lambda_{j,t}(p)-\mathrm{lo}(p)\bigr)-m\,\mathrm{fl}(p)
\ \le\ \sum_{j=1}^{m}\mathrm{gain}(j,w)-\tau(w),
\end{equation}
where $\Lambda_{j,t}(p)$ are von Mangoldt class sums at $p$ and $\mathrm{lo}(p)$, $\mathrm{fl}(p)$ explicit error
terms (with $b^{\delta(x)}$ in place of $b$); and (iv) a containment of served orbits and a bound for the mean of
$\kappa$. Here $a$ is a size parameter of $L$ and $TT$ a finite set attached to the host. Inequality
\eqref{eq:families:linkG} says that the divisibility gains at $w$, summed over the $m$ lifts, dominate the
von Mangoldt class sums of the prime $p$ up to the error terms.

\begin{theorem}[Local lattice link]\Fmark\label{thm:linkG}
Let $\mathcal C$ be a family class on which the group compatibility of the extra data and the condition
$\sfname{Adm2F}$ (eventually $n\delta_1/2\le|TT|\le2n\delta_1$ and $B_{TT}\le|TT|\,x^{-2+\eps}$) are closed. Then
$\sfname{LinkG}$ is closed on $\mathcal C$.
\formal{Formalism.L2.G.linkG_concrete5}
\end{theorem}

\begin{proofoutline}
On $D3(3,\mathcal A)$ both hypotheses hold for every member: $\sfname{Adm2F}$ by a Mertens-type estimate
\cite{Mertens1874}, group compatibility from the duties and the selection. Group compatibility gives trivial
stabilisers on finite classes and a standard-subgroup property; ten closed properties are then combined, the main
one being an orbit-valuation bound with the absolute constant of \cref{thm:families:local}. An order estimate, a
Weyl-sandwich condition, orbit-lattice bounds (a linear-programming input) and a lattice sandwich then give the
inequalities of $\sfname{LinkG}$ (about $80\,000$ lines).

%P2PTR thm:linkG
A complete proof is given in \cref{sec:pf:linkG}.
\end{proofoutline}

The deepest ingredient is a purely local statement about lattices in algebraic representations. An
\emph{unramified quadratic setting} consists of the unramified extension $W$ of $\Z_p$ of degree $n'$ (even) with
Frobenius $\sigma$, and a finite extension $\OO\supseteq W$ with uniformiser $w$, $w^2=p$, $|\OO/w|=p^{2n'}$, in
which every rational integer is a square. Lattices $L\subset W[1/p]^{m+1}$ are joined by \emph{Hecke edges of type
$r$}: pairs $e_s\subseteq e_t$ with $p\,e_t\subseteq e_s$ and $[e_t:e_s]=p^{n'r}$. For a weight
$\lambda_0:\Z/n'\to\Z^{m+1}$, the $\sigma$-twisted tensor product $V$ of dual Weyl modules contains weight
lattices $H_L$ attached to lattices $L$ (intersections of the integral structures
$\det(g)^{-\lambda_m}\nabla(\lambda-\lambda_m,g)$ over bases $g$ of $L$, cf.\ \cite{Jantzen2003}), lattices $R_e$
attached to edges, and a perfect pairing $b$ with the dual. $\sfname{LPprime}(t,D,\kappa,c_*)$ asserts:
$4(D+1)\le c_*$, and for \emph{every} lattice $L$ there are edges $e_1,\dots,e_M$ with target $L$ and vectors
$F_i\in w^{-t}H_{e_{i,s}}\cap\bigcap_{e:\,e_s=e_{i,s}}R_e$ such
that $F_i\in w^{c_*}R_{e'}$ for every other edge $e'$ with the same source, and the pairings $G\mapsto b(F_i,G)$
have, on the dual weight lattice at $L$ and up to the denominator $p^D$, rank at least $(1-\kappa)\dim V$.

\begin{theorem}[Local lattice theorem]\Fmark\label{thm:families:local}
There is an absolute constant $C_0\ge0$ with the following property. Let an unramified quadratic setting with
$p\ge13$ be given, integers $r\ge2$ and $m$, a set $TT\subseteq\Z/n'$, a dominant weight $\lambda_0$ adapted to
$TT$ with entries bounded in terms of parameters $a,c_{pp},c_0$, and numbers $L_1,L_2,D_p,c_*,Y$; let $N_1$ and
$G'$ be the explicit orbit-window counts attached to auxiliary reals $(\delta,h,x)$. Assume the orbit hypotheses (a
tiling condition for $TT$, $2L_1\le L_2<\lceil1/\delta\rceil$, $3G'\le n'$), $N_1^2\le r|TT|$, $r\le Y$,
$C_0N_1^2\log a\le m+1$, $N_1+2r\le m+1-r$, $256(m+1)^2|TT|Y^2\le p^{L_1-1}$, $C_0|TT|(m+1)^2\log_pa\le D_p$,
$8(D_p+1)\le c_*$ and
\[
C_0\,(m+1)^3r^2|TT|\,(p^{G'})^5\bigl(c_*+r|TT|c_0+2rn'c_{pp}\bigr)\,Y\ \le\ a .
\]
Then $\sfname{LPprime}\bigl(t,D_p,C_0|TT|r^2(m+1)/Y,c_*\bigr)$ holds for every $t$ beyond an explicit bound.
\formal{Formalism.MainNotions.G.second_theorem_absolute_constant}
\end{theorem}

\begin{proofoutline}
A criterion reduces the conclusion to thirty-seven explicit clauses about tensor-product operators over the
Frobenius orbit $\Z/n'$, operator weight bounds and a numerical counting statement (about $31\,000$ lines); the
clauses are verified using extremal projectors, a plane count, restriction estimates and path weights (about
$15\,000$ lines), Fock-model lattice lemmas, and the Weyl dimension formula, branching rule and Pieri rule for
Gelfand--Tsetlin patterns \cite{Weyl1946,GelfandTsetlin1950}. In words: from every lattice one reaches, along Hecke
edges, vectors of the weight lattices which are highly divisible on all competing edges and whose pairings have
almost full rank.

%P2PTR thm:families:local
A complete proof is given in \cref{sec:pf:linkG}.
\end{proofoutline}

\paragraph{Division, orthogonality and the remaining links.}
$\sfname{LinkDIV}$ collects, eventually and for every word $u$: a ``one map'' condition; the frame bound
$s_{\mathrm{Frame}}(N_u,h)+\mu_{\max}(\mathit{RR})\le x^{C_H}b$; the \emph{Bessel identity} that the functionals
$B_u$ agree, up to non-zero scalars and isometries, with one fixed pairing; a height bound
$\mathrm{ht}(B_k(u))\le h(u)$; an upper bound for $\mu_{\mathrm{dual}}+h(u)$ by a main term plus errors; and two
approximate identities relating $h(u)$, the level primes $T_u$ and the type primes. $\sfname{LinkDIVt}$ is the
implication $\sfname{LinkEta}\wedge\sfname{DIVinN}\Rightarrow\sfname{LinkDIV}$.

\begin{proposition}[Division link]\Fmark\label{prop:families:div}
If, for kernel data attached to the members of a class $\mathcal D$, the ``division kernel constant bounds'' are
closed on $\mathcal D$ under $\sfname{TuSize}$ and $\sfname{DGamSize}$, then so is $\sfname{LinkDIVt}$. On
$\familiesDiii$ these bounds follow from the fields of the package: a frame-number bound (level wiring, a slope
lemma, a Mertens-type estimate for $\tau$-sums and a bound for $\mu_{\max}$ under $\sfname{TuSize}$); the
integrality of period pairings at all places, from the period rows and the service row
(\cref{prop:families:periods}); stability of the lattices under $\OO_{\mathit{LL}}$; complement conditions; and
rationality of the Bessel functionals.
\formal{Formalism.L2.DIV.closed_linkDIVt_of_closed_bridge_kernel}
\end{proposition}
%P2PTR prop:families:div
\begin{proof}See \cref{sec:pf:top-shape}.
\end{proof}


The Bessel identity reflects the seesaw step of \cite{Announcement2026}, in which the period of the iterated lift
is identified with a Bessel period independent of the Hecke characters (\cref{rem:families:seesaw}).

\begin{corollary}[Approximate orthogonality link]\Fmark\label{cor:families:linkS}
Under the spectral selection and the fields $\mathsf{hWu}$, $\mathsf{hRF}$, $\mathsf{hTF}$, $\mathsf{c30}$,
$\mathsf{sTy}$, $\mathsf{sNG}$, $\mathsf{sM}$, $\mathsf{sLp}$, $\mathsf{sO}$, $\mathsf{h30}$, $\mathsf{hInstR}$ of
the package, $\sfname{TuSize}\Rightarrow\sfname{LinkS}(3)$ is closed on $\familiesDiii$. Here
$\sfname{LinkS}(C_0)$ asserts that eventually there are $\eta_0,\eta_1,e_0\ge0$ with
$\eta_0,\eta_1,\sqrt{e_0}\le\exp(-x^{\eps/8})$ such that, with $\Gamma=\exp(C_\Gamma m^5f\Lambda)$, if
$C_0^{-1}\le d_u\Gamma^2/D_*\le C_0$ for all words $u$, then
\[
\mathrm{hs}(u,u')\le(1+\eta_0)\frac{d_ud_{u'}}{D_*}\ (u\ne u'),\qquad
\|K_{uu}-\mathrm{id}\|_{\mathrm{HS}}^2\le e_0\,d_u,\qquad D_{\mathrm{all}}\le(1+\eta_1)D_*,
\]
where $\mathrm{hs}(u,u')$ is the Hilbert--Schmidt mass of the off-diagonal block.
\formal{GlobalTheory.Top.link_alt_alt_bis}
\end{corollary}

\begin{proof}[Proof sketch]
The selection provides local ties and an instance tie; residues are eventually controlled; a transport form of
\cref{thm:theorem-S} turns the row estimates (type, non-generic and main-term rows), the loop rows and the point
count into the three Hilbert--Schmidt estimates.

%P2PTR cor:families:linkS
A complete proof is given in \cref{sec:pf:top-shape}.
\end{proof}

\begin{remark}\label{rem:families:small}
$\sfname{LinkEta}$ (the local places of $D$ cover its twisted places) holds for every member of $\familiesDiii$,
by the normalized local data of the tied records of the twist selection. $\sfname{SmallLevelRankEq}$
($d_{Z_0}(j)=D_0$ for all $j$) follows from the primitivity, $\lambda$- and division rows together with the
surjectivity of the level projections. $\sfname{LinkMu}$,
$|\mu(N_u)-\mu(N^H_u)-\sum_w\mathrm{sz}(w,\sigma(u,w))|\le C\,X''(n)$, is deduced from $\sfname{LinkDIV}$ in the
assembly of \cref{sec:stepanov}.
\end{remark}

\paragraph{Mass ratios and slope inequalities.}
The global theorem (\cref{thm:global-chain}) needs eleven closed properties of the class, collected in a record
$\sfname{FamGLmin}$. For a word $u$ let $\Lambda_u$ be the hermitian lattice at level $0$ attached to $u$; for
embedding types $\tau$ let $a_\tau,a'_\tau$ be column sizes, $s_\tau$ the spread of the first column and
$\lambda_\tau$ its last entry.

\begin{proposition}[Mass ratios and slopes]\Fmark\label{prop:families:famgl}
Let $\mathcal A$ satisfy the duties, and assume the field $\mathsf{tie}$, the volume ratio $\mathsf{hU}$ and the
rank-two genus mass $\mathsf{hG}$. Then, closedly on $\familiesDiii$: the mass ratios
$\mathrm{mass}(D,M)/\mathrm{mass}(D,M-1)$ ($M\ge3$) are those of \cref{prop:families:mass}; the base volume equals
the mass in rank $r$ (this is deduced from the rank-two genus mass $\mathsf{hG}$ and the volume ratio $\mathsf{hU}$);
the $\sigma$-degrees are bounded on both sides in terms of the gaps; for every word $u$ with
column valuations,
\begin{gather*}
\mu(\Lambda_u)\ \ge\ -\sum_\tau\frac{a_\tau}{k}\log T_u-\sum_\tau\bigl(a'_\tau\log k+k^2\log(s_\tau+k)\bigr),\\
\mu_{\max}(\Lambda_u)\ \le\ \sum_\tau\bigl(4a'_\tau\log(kv)+k^2\log(s_\tau+k)+|\lambda_\tau|\log T_u\bigr),
\end{gather*}
with $v$ a constant of the host; and column valuations, free $T$-data, slope lines, Hecke integrality at level
zero and a second-branch inclusion length hold. Consequently $\sfname{ChosenE}\Rightarrow
\sfname{NuFullF}\wedge\sfname{MuRank}$ and $\sfname{TuSize}\Rightarrow\sfname{DGamSize}\Rightarrow
\sfname{LinkGLOBAL}$ are closed on $\familiesDiii$.
\formal{GlobalTheory.Top.link_gl}
\end{proposition}

\begin{proofoutline}
The slope inequalities use exact slopes of column lattices, the archimedean class map, a short-dual-basis
lemma, a base-change formula for Hermite--Siegel slopes and a hook-length bound; they are Arakelov slope
estimates of the kind used by Bost \cite{Bost1996,Bost2001}. The proof also uses that every totally positive
element of a CM field unramified over its real subfield is of the form $as\bar s+bt\bar t$ for given totally
positive $a,b$. The last assertion is \cref{thm:global-chain} applied to the eleven properties.

%P2PTR prop:families:famgl
A complete proof is given in \cref{sec:pf:top-shape}.
\end{proofoutline}

\paragraph{Designed families.}
\begin{proposition}[Existence of designed families]\Fmark\label{prop:families:design}
Let $\mathcal A$ be a class with a twist selection and the duties, satisfying $\mathsf{hTu}$, $\mathsf{hn4}$ and
$\sfname{LitThin}$, and such that $\sfname{MinkDiv}$ implies design choices for $\sfname{WithTw}(\mathcal A)$ (in
the package: the field $\mathsf{hEx}$, applied for $\lambda>1$ and a requirement $G$ (\cref{def:requirement}) with
$\sfname{HostExW}(\lambda,G)$ to the class $\mathcal A\wedge G.\mathrm{cls}_\lambda$). Then designed families exist: for
small $\eps$ and every staged threshold there are large families in $D3(3,\mathcal A)$.
\formal{GlobalTheory.Top.designExists_of_tuSize_aux4n_litThin}
\end{proposition}

\begin{proofoutline}
$\sfname{MinkDiv}$, a Minkowski-type statement on multipliers, is proved in the development, so the hypothesis
yields design choices; they become designed families using $\mathsf{hTu}$, $\mathsf{hn4}$ and a ``thin'' choice of
auxiliary primes in short intervals derived from $\sfname{LitThin}$ (classical zero-free regions with the
Landau--Page exceptional zero, the truncated explicit formula and zero-density estimates near $\Re s=1$ for
Dirichlet $L$-functions \cite{Davenport2000,Montgomery1971}, all proved in the development).

%P2PTR prop:families:design
A complete proof is given in \cref{sec:pf:top-shape}.
\end{proofoutline}

% ---------------------------------------------------------------------------
\subsection{Point inequalities, member bundles and non-emptiness}\label{sec:families:nonempty}

\paragraph{The four point inequalities.}
A point of a family over $(x,i)$ carries two \emph{indefinite data} $S_{\mathrm f},S_{\mathrm t}$ (the two sides
of its complex realisation) and numbers $N(C,C_H)$. An indefinite datum $S$ consists of a set $\mathrm{Tw}$ of
twisted places, an integer $D(S)$, lattices $M_j$ with hermitian metrics $X_{W,j}$ at the infinite places, maps
$\theta_{j,t}$ of theta type sending lattices $T_{j,t}$ into the spaces of the $M_j$ and multiplying hermitian
majorants by scalars $\lambda_{j,t,p,v}>0$, scalars $\kappa_v$, error terms $\mathrm{err}_0(w,j)$, and the glued
lattice $\mathit{RR}$; $I_j(t,w)$ is the $w$-adic generalised index of $\theta_{j,t}(T_{j,t})$ in $M_j$. The numbers
are $l_3(j)=l_{3m}(j)+Cl_{3u}$, $l_6(j)$, $E_1(j)=CE_{1u}$ and $hh=x^{C_H}b$. The four inequalities are:
\begin{itemize}
\item $\sfname{L3}$ (covolume): $\log\covol(M_j,X_{W,j})/\rank\mathit{RR}\le l_3(j)$ for all $j$;
\item $\sfname{ON}$ (almost surjectivity): $\mathrm{err}_0\ge0$, and at every twisted place $w$ (residue
characteristic $p_w$, ramification index $e_w$, $\ell_w=\log\Nm w/[\mathit{LL}:\Q]$)
\[
\sum_j\bigl(I_j(\mu_w,w)+I_j(\mu_w,\bar w)\bigr)\le D(S)\Bigl(\sum_j e_w\bigl(2F_0(w,j)+\mathrm{err}_0(w,j)\bigr)
+\frac{2e_w}{p_w^2}\Bigr),
\]
and $\sum_{w\in\mathrm{Tw}}\ell_w e_w\bigl(\mathrm{err}_0(w,j)+J_0(w,j)\bigr)\le E_1(j)$;
\item $\sfname{L6}$ (low weight): $l_6(j)\le\log\lambda_{j,\mu_0,p,v}$ for all $j,p,v$;
\item $\sfname{FUN}$ (slope): $\mu_{\max}(\mathit{RR})\le hh-\mathrm{meanLog}(\kappa)$.
\end{itemize}
The family statements are: $\sfname{CovolumeBoundOfX}$ (closedly on a class $\sfname{FamHyp}$, there is $C_0$ such
that eventually at all points $S_{\mathrm f}$ satisfies $\sfname{L3}$ for all $C\ge C_0$), proved in
\cref{sec:covolume}; $\sfname{SurjectiveOntoX}$ (the same for $\sfname{ON}$, under $\sfname{DIVinN}$);
$\sfname{LambdaLowerBound}_2$ ($S_{\mathrm t}$ satisfies $\sfname{L6}$ for all $C$); and
$\sfname{ConstantFunctionFits}$ ($S_{\mathrm f}$ satisfies $\sfname{FUN}$ at points carrying a grid at level
$C_F$, for all large $C_F$).

\begin{proposition}[Surjectivity, low weight and slopes]\Fmark\label{prop:families:points}
$P$ implies $\sfname{SurjectiveOntoX}$ and $\sfname{LambdaLowerBound}_2$; and $F$, two components of $P$ (a join
statement and the functional equation of Hecke $L$-functions), \cref{thm:theorem-K}, $\sfname{FunctionalsSeparate}$
and $\sfname{LambdaLowerBound}_2$ imply $\sfname{ConstantFunctionFits}$.
\formal{GlobalTheory.Fc.S0.ontoX_and_lamLow_of_unitaryGroupHypotheses'}
\formal{GlobalTheory.BigTop.Fun.Close.fourier_coefficient_fit_of_join_alt}
\end{proposition}

\begin{proofoutline}
$\sfname{ON}$ reduces to a bound for the weighted index excess at the twisted places, proved through ``record
objects'' attached to points (twist bounds at conjugate pairs of places, share bounds for eigenspace dimensions,
depth bounds at good, ramified and level primes, and an exponent bound $\sum_\ell J_p(j,\ell)\log\ell\le E_J$) from
Green's theorem, Clozel's purity, $\ell$-adic Hecke characters, Chebotarev, Galois representations of polarized
cuspidal representations, base change, the crystalline lattices of Bhatt--Scholze and archimedean theory (about
$280\,000$ lines). $\sfname{L6}$ comes from an inner product formula: the lift composed with its adjoint acts on
each type part by a scalar $l>0$ with $\log l\ge\mathrm{mainTerm}-8fn_V^2L_p$, proved from the Rallis inner product
formula for the twisted lift and the identification of the finite zeta integral with a partial $L$-value.
$\sfname{FUN}$ is an Arakelov bound for the maximal slope: the Petersson norm of the bottom lift is a Rallis
integral \cite{Rallis1984,KudlaRallis1994}, the class sums are bounded at all places, and a numerical slope
inequality concludes.

%P2PTR prop:families:points
A complete proof is given in \cref{sec:pf:onto3,sec:pf:fit}.
\end{proofoutline}

The definite-side condition $\sfname{DefiniteSideCondition}(b,C,C_H)$ of a bundle $b$ is a height bound under the
point inequalities: given witnesses of the primitive-row conditions, of $\sfname{L3}$, $\sfname{ON}$, $\sfname{L6}$
and of two isometry conditions, every word $u$ admits a column chain, a frame $X_0$ at the infinite places with
$(X_0^{-1})^*X_0^{-1}$ equal to the hermitian metric, and height numbers bounding the operator norm of the final
point map by $e^{h/2}$ at every infinite place. The field $\mathsf{rRows}$ asserts it along the selection.

\paragraph{Twisting data and members.}
\begin{definition}[Twisting data; large unramified exponents]\label{def:families:twist}
A \emph{twisting datum} over a base is a finite set $T$ of rational primes together with, for each place $w$ of the
coefficient field with residue characteristic $\ell_w$, a finite set $Q_w$ of primes of $\OO_E$ not above
$\ell_w$, an exponent $k_w\in\N$ and a Hecke character $\psi_w$ of $E$ of order dividing $\ell_w^{k_w}$, trivial on
the id\`eles of $F$, tame, and ramified only at split primes of prime norm; $\psi_w=1$ and $k_w=0$ if
$\ell_w\notin T$. A property holds \emph{eventually for large twisting data} ($\sfname{EvLgG}$) if, eventually in
the scale, there are finite exceptional sets of algebraic numbers and of primes, depending on the base and the
frame, such that the property holds for all bases whose coefficient fields contain these numbers and all twisting
data whose prime set contains the prescribed primes and whose sets $Q_w$ avoid the excluded ones. (This is the
meaning of ``large unramified exponents'' in the names of the development.)
\end{definition}

For a parameter $p$, $\sfname{MemberE}(p)$ asserts: bases with the required coefficient fields exist; for large
untwisted data a built point exists; built points satisfy export conditions at the twisted places; and every
built point can be moved, by enlarging the base and re-twisting with the same prime set, to a point which is
built, good and twisted at all places. $\sfname{ExBL}(p)$ asserts that there is a slot $R$ containing, eventually
and for large twisting data, all built points, such that every good point of $R$ carries linked bundles of every
size. The parameters used are $p_{\mathrm{dep}}(h_L,c)=\sfname{parametersOfDependence}(h_L,c)$, whose norm-choice
extension asks for a law (rational map values are rational) and a grid datum at level $(1,c,b)$ with full
$\lambda$-tie, and $p_A(h_L,c)$, which adds a saturated state of a ``strict-henselian acyclic'' level and an
at-property $A$; the norm-choice cores contain local data built from $h_L:\sfname{RecTExists}$.

\begin{theorem}[Member bundles]\Fmark\label{thm:member-bundles}
Assume $P$ and $\sfname{RowRat}$. Let $h_L:\sfname{RecTExists}$, $c\in\R$, and let $A$ be an at-property with
$\sfname{EvLgG}(A)$. Then the Siegel mass ratio of \cref{prop:families:mass} holds, and $\sfname{MemberE}(p_A)$
and $\sfname{ExBL}(p_A)$ hold for $p_A=p_A(h_L,c)$.
\formal{GlobalTheory.T.large_unramified_exponent_bundles_exist_alt}
\end{theorem}

\begin{proofoutline}
From the fields of $P$ and $\sfname{RowRat}$ one obtains saturated states without localisation primes and equality
of twisted multiplicities at every weighted point (\cref{rem:families:mult}). Member existence runs through a
chain of levels of a state (same-Galois, intermediate levels with rationality and cuspidality of the rows, the
strict-henselian acyclic level with its height-counting character data, and levels three and four with norm
choices, theta data, column data at split primes and real ties; about $207\,000$ lines). Export realises the twist
series demanded at twisted places from Galois cohomology and the Galois representations of polarized cuspidal
representations. Bundles are produced by a recursion on an $\eta$-size, using uniform zeta data and a reduction to
local Langlands parameters; the at-property is added using $\sfname{EvLgG}(A)$.

%P2PTR thm:member-bundles
A complete proof is given in \cref{sec:pf:bundles}.
\end{proofoutline}

\begin{theorem}[Open parameter statements]\Fmark\label{thm:families:open}
$P$ implies, for every $h_L$ and $c$, twelve linked statements for $p_{\mathrm{dep}}(h_L,c)$: primitivity (Hecke
rationality, unit and equivariance conditions, block condition, $\sfname{RowTest}$, separated types),
$\lambda$-rows, division rows, decomposition points, existence of a join and of a comparison, labels (a Hida
setting, a good operator and an isometry with label ties), goodness, geometric compatibility, a multiplicity
condition, a $\kappa$-condition, and eventually $\sfname{DefiniteSideCondition}$ for all $C\ge C_0$.
\formal{GlobalTheory.T.open0_of_unitary_group_hypotheses}
\end{theorem}
%P2PTR thm:families:open
\begin{proof}See \cref{sec:pf:open0}.
\end{proof}


In the development these are called ``open orbit'' statements (stage $0$ of the row construction); no
prehomogeneous vector space is involved.

\paragraph{The main theorem.}
\begin{theorem}[Non-emptiness of the package]\Fmark\label{thm:hyps-nonempty}
Let $P$ and $F$ be the two bundles of classical results of \cref{sec:families:theta}, and let $h_A$, $h_M$,
$h_X$, $h_L$ be proofs of $\sfname{ThetaIntegralsAgree}$, $\sfname{LocalMatching}$, $\sfname{CovolumeBoundOfX}$
and $\sfname{RecTExists}$. Assume that for every $c\in\R$
\[
\familiesLinked\bigl(p_{\mathrm{dep}}(h_L,c),\ \familiesRow(\cdot,\cdot,\mathcal A\mapsto
\familiesClosed(D3(3,\mathcal A),\sfname{BoundAtTwoThree}))\bigr).
\]
Then the package $\sfname{HypsAlt}(\sfname{NL},\sfname{DefiniteSideCondition})$ is non-empty.
\formal{UpperTheory.Lit.nonempty_hyps_definite_side_of_closed_bound}
\end{theorem}

\begin{proofoutline}
The proof assembles the results above.
\begin{enumerate}[label=\arabic*.]
\item \cref{prop:families:fs} gives $\sfname{FunctionalsSeparate}$, and \cref{prop:families:rowrat} gives
$\sfname{RowRat}$.
\item \cref{prop:families:points} gives $\sfname{ConstantFunctionFits}$. From $F$, $h_A$, $h_M$ and
$\sfname{FunctionalsSeparate}$ one obtains a constant $C_{F1}$ such that full grids exist eventually for large
twisting data at every level $C_F\ge C_{F1}$, and from the fit a constant $C_{F0}$ for the functional rows. Put
$c^*=\max(C_{F0},C_{F1})$ and $A=\sfname{GridA}(c^*)$, so that $\sfname{EvLgG}(A)$.
\item \cref{thm:families:open} gives the open statements for $p_{\mathrm{dep}}(h_L,c^*)$; by monotonicity they hold
for $p_A(h_L,c^*)$, which yields the functional row ($\sfname{FUN}$) and, with $h_X$, the covolume row
($\sfname{L3}$) for $p_A$.
\item \cref{thm:member-bundles} gives the mass ratio, $\sfname{MemberE}(p_A)$ and $\sfname{ExBL}(p_A)$.
\cref{thm:linkN}(b), applied to the hypothesis at $c=c^*$, gives the norm row for $p_{\mathrm{dep}}$, hence for
$p_A$; \cref{prop:families:allrows} gives $\sfname{AllRowsAlt}$ for $p_A$.
\item \emph{Class selection.} The twenty-six linked rows and a size row from the genus mass formula are combined
into one linked row. A class-selection theorem turns it, together with $\sfname{MemberE}$ and $\sfname{ExBL}$, into
a class $\mathcal A$ of families built from good points and a selection $\sigma$ of linked bundles satisfying
all rows; $\sfname{MemberE}$ gives the design choices of $\mathsf{hEx}$, and the remaining fields are uniform
theorems of the development.\qedhere
\end{enumerate}

%P2PTR thm:hyps-nonempty
A complete proof is given in \cref{sec:pf:stageU}.
\end{proofoutline}

\begin{remark}[Supply of the hypotheses]\label{rem:families:supply}
In \cref{sec:top} the theorem is applied with $P$, $F$ the structures of classical results proved in the
development, $h_A,h_M$ from \cref{thm:theorem-K}, $h_X$ from \cref{sec:covolume}, $h_L$ from \cref{sec:galois}, and
the linkedness hypothesis from \cref{lem:families:at23}; thus the package is non-empty unconditionally
(see the proof of \cref{thm:zerofreeall}). The covolume bound enters only through the covolume row, hence only through
$\mathsf{rRows}$; $h_L$ enters through the norm-choice data of the points, the Galois preconditions
(\cref{thm:linkN}) and the open statements. Apart from $\mathsf{hEx}$, the fields are closed statements,
statements about members or selected bundles, or uniform statements, and a closed statement holds trivially for
a class without large members; the existence of large members is the content of $\mathsf{hEx}$, which comes from
\cref{thm:member-bundles} through class selection.
\end{remark}

{\small
\begin{longtable}{>{\raggedright\arraybackslash}p{0.36\textwidth}>{\raggedright\arraybackslash}p{0.58\textwidth}}
\caption{Fields of $\sfname{HypsAlt}$ and their sources.}\label{tab:families:fields}\\
\toprule
fields & source\\
\midrule
\endfirsthead
\toprule
fields & source\\
\midrule
\endhead
$\mathcal A$, $\sigma$, $\mathsf{hEx}$, $\mathsf{hduty}$, $\mathsf{hMt}$ & class selection from
\cref{thm:member-bundles}\\
$\mathsf{hPA}$, $\mathsf{hLamP}$, $\mathsf{hDivG}$, $\mathsf{hDcPtG}$, $\mathsf{selJ}$; $\mathsf{jOn}$, $\mathsf{zL}$,
$\mathsf{hTw}$ & \cref{thm:families:open}, directly or by short deductions\\
$\mathsf{hRB}$, $\mathsf{hn4}$, $\mathsf{hInstR}$, $\mathsf{aC2}$, $\mathsf{aAtG}$ & archimedean inputs, class-number
formula, volume ratio (\cref{prop:families:volratio}), conditions at constructed points\\
$\mathsf{aPt0}$, $\mathsf{aPtTw}$, $\mathsf{aPer0}$, $\mathsf{aPerTw}$, $\mathsf{aP10}$, $\mathsf{tie}$ &
\cref{prop:families:periods}\\
$\mathsf{sLp}$, $\mathsf{sO}$, $\mathsf{hTF}$; $\mathsf{hServ}$ & loop rows, \cref{thm:theorem-S}, trace frame
(Siegel--Weil); CM theory from Artin reciprocity (\cref{sec:families:rows})\\
$\mathsf{c30}$, $\mathsf{hc30}$, $\mathsf{h30}$, $\mathsf{cIc}$, $\mathsf{hcI}$, $\mathsf{irrC}$ & \cref{thm:counting}
and the label statement\\
$\mathsf{rRows}$; $\mathsf{nRest}$ & open statements, \cref{prop:families:points} and $h_X$; \cref{thm:linkN}\\
$\mathsf{cT}$, $\mathsf{hTu}$, $\mathsf{hLit}$, $\mathsf{hWu}$, $\mathsf{hRF}$, $\mathsf{sTy}$, $\mathsf{sNG}$,
$\mathsf{sM}$, $\mathsf{hU}$, $\mathsf{hG}$ & an explicit threshold function; uniform statements proved in the
development\\
\bottomrule
\end{longtable}
}
